\documentclass[10pt]{article}
\usepackage[letterpaper,margin=0.6in,top=0.55in,bottom=0.6in]{geometry}
\usepackage[T1]{fontenc}
\usepackage{lmodern,microtype,enumitem,fancyhdr}
\renewcommand{\familydefault}{\sfdefault}
\setlist{nosep,leftmargin=1.2em,itemsep=1.5pt}
\setlength{\parindent}{0pt}\setlength{\parskip}{4pt}
\pagestyle{fancy}\fancyhf{}\renewcommand{\headrulewidth}{0pt}
\fancyfoot[L]{\footnotesize Bellingham Math Circle / Week 10 / Adult table card / Unpiloted}
\fancyfoot[R]{\footnotesize Full answers: the facilitator guide}
\newcommand{\sect}[1]{\par\medskip{\large\bfseries #1}\par\smallskip}
\begin{document}
{\footnotesize Week 10 / Bridges / For the adult at each table}\par
{\LARGE\bfseries Bridges: the one-page version}\par
\sect{What the activity is}
Each drawing is a town: white circles are islands, grey bands are bridges. A counter sits on every bridge; a token walks from island to island, and crossing a bridge means picking up its counter, so no bridge can be used twice. The question all week: can you pick up every counter, and which island do you have to start on?

\sect{Materials (whole room; about 20 minutes to gather and test)}
\begin{itemize}
\item Counters about 1 inch or smaller: about 20 per pair at K--1, 26 per pair at Grades 2--3, 18 at Grades 4--5 (about 150 in all). Tokens that look different from the counters, 2 per pair (one walks, one marks the start).
\item Yellow pattern-block hexagons (about 40) and 4.5-inch craft sticks (about 60) for the building problems (K--1 P7, Grades 2--3 P8). Lay each stick so it touches both hexagons; its counter sits on the stick. Two sticks side by side are two bridges between the same pair.
\item Coloured pencils, pencils, scrap paper; small sticky notes to letter hexagons (Grades 2--3). Masking tape for the optional floor town.
\item Print single-sided at 100\%, not ``fit to page'', or the bridges shrink below counter size. Before the session, check that counters fit the tightest town (K--1 page 9, second town).
\end{itemize}

\sect{Launch (two minutes, whole group)}
Use a house-shaped town (five islands, six bridges) taped on the floor or K--1 page 1. ``Every bridge has a counter. Crossing a bridge, you pick up its counter; a bridge with no counter is closed.'' A child starts on a bottom corner, walks, and gets stuck with counters left. ``Would a different start pick up every counter?'' Don't say which start works.

\sect{How it runs, and what Problem 1 is for}
\begin{itemize}
\item \textbf{Grouping:} K--1 and Grades 2--3 in pairs (one walks, one keeps the counters and checks); Grades 4--5 a pair plus a referee who checks every move.\item \textbf{Every band:} put every counter back before each new try. Two counters stacked on one bridge means it can be crossed twice (take the top one). Each tracing picture is printed twice for a second try.
\item \textbf{Problem 1} is a few small towns to walk: K--1 colours the island it started on, and Grades 2--3 and 4--5 write the walk as island letters, like A B C A (show this; the 4--5 page gives no example). Its job is to get the walking rule into hands; children find that some starts work and some get stuck.
\item Core of the hour: K--1 Problems 1--3, Grades 2--3 Problems 1--4, Grades 4--5 Problems 1--5. The rest is for children who get further.
\end{itemize}

\sect{The surprise (don't tell it; let the counters show it)}
Count the bridges at each island. If every island has an even number, start anywhere and you end where you started. If exactly two islands are odd, start at one and you end at the other. If more are odd, it can't be done: every time you pass through an island you use its bridges two at a time, so an odd island must be the start or the end.

\sect{Key answers}
\begin{itemize}
\item \textbf{K--1.} P1: every town can be done; the house needs a bottom corner (an island with 3 bridges), and a town whose islands all have an even number works from anywhere. P2: two squares, only the two middle islands; big triangle, any island. P3: X, yes, yes, yes, X (the three-arm star fails: three tips, each with one bridge). P4: yes, X, X, yes, yes. P5: cross out the rectangle with an X in it and the circle with a plus (the house with an X can be traced). P6: join any two islands that have 1 or 3 bridges. P8 (pages 11--12): arc town, all 6 bridges; the town with two tails, the 2 tail bridges; square with a centre island, the 4 sides.
\item \textbf{Grades 2--3} (Grades 4--5 P1 and P2 use the same towns). P1: house from B or C; triangle with tail from A or C; two double bridges from anywhere. P2: town 1 none, town 2 A or D, town 3 D or F, town 4 anywhere. P3 comes before the rule, so let children guess, then check with counters: big triangle, all six; triangle with a double bridge, A and C; arc town and square with three spokes, none. P4: join any two islands with an odd number of bridges. P6 is the rule above. P9: impossible; walk it backwards, or round its loop from another island, to get a second start.
\item \textbf{Grades 4--5.} P1--2 as Grades 2--3. P3: 2 by 4 ladder, none; bowtie, anywhere; 3 by 3 grid, F to H or H to F. P4 is the rule and why. P5, K\"onigsberg: impossible, since all four land areas are odd; one new bridge between any two areas makes a walk possible, and two new bridges make one that ends where it started. P6: no and no; the bridge counts add up to twice the number of bridges, an even number, so the odd islands always come in pairs.
\end{itemize}
\end{document}
